\section{Multiscale proposal models and audited distribution shift}
\label{sec:multiscale-extension}
This extension separates three mathematical objects: the unknown attack law,
the fitted proposal function, and a confidence bound from new attack trials.
Its results do not replace the registered sequential gate. In particular,
a derivative of a neural network describes its response to an input change;
it does not bound an adversary's true success probability. The new software is
\texttt{privacy\_risk.multiscale}. All new numerical demonstrations use synthetic
laws and do not modify the earlier public-data results.

\subsection{An arbitrary dimensional risk vector}
Let $x\in\mathcal X=\prod_{j=1}^d[\ell_j,u_j]$ describe a release and its
registered threat model, with $u_j>\ell_j$. For coordinate $i\in\{1,\ldots,m\}$,
let $A_i(Z,x)\in\{0,1\}$ be the declared attack success event. Define
\begin{equation}
 r_i(x;P_i)=\int A_i(z,x)\,dP_i(z),\qquad
 z_i(x)=1-r_i(x;P_i)/\tau_i,\quad 0<\tau_i<1.
\end{equation}
A coordinate includes its target population, information available to the
attacker, attack algorithm and harm. These components cannot change silently
when the feature vector changes. Eligibility is the intersection of all
registered inequalities $r_i\le\tau_i$, rather than an average over harms.
The scalar $\min_i z_i$ records the limiting margin, but a report retains the
whole vector so that a favourable component cannot conceal another violation.

\subsection{Finite functional wavelet dictionary}
Normalize $v_j=(x_j-\ell_j)/(u_j-\ell_j)$ and define
\begin{equation}
 \psi(t)=(1-t^2)e^{-t^2/2},\qquad
 \psi'(t)=(t^3-3t)e^{-t^2/2}.
\end{equation}
For levels $a=0,\ldots,J-1$ and centres
$c_{ak}=(k+1/2)2^{-a}$, $k=0,\ldots,2^a-1$, the feature map is
\begin{equation}
 \Phi_J(x)=\big(v_j,\psi(2^a(v_j-c_{ak}))\big)_{j,a,k}
 \in\mathbb R^{D},\qquad D=d\,2^J.
\end{equation}
The dictionary captures localized changes at several scales and retains the
normalized identity coordinates. It is neither an orthogonal wavelet transform
nor an exact interpolant. On a bounded interval the atoms need not be orthogonal.
An exact interpolation claim would require specified nodes, an invertible
collocation matrix and a conditioning estimate, none of which follows from
using a wavelet-shaped activation. Tensor products of all atoms cost
exponentially in $d$; the coordinatewise dictionary avoids that construction.
Cross-feature interactions arise in the dense network. Its computational benefit
is a choice of representation, not a universal approximation-error guarantee.

Put $s_j=(u_j-\ell_j)^{-1}$. Each feature depends on a single raw input, so the
feature Jacobian has orthogonal columns. Since
$|\psi'|\le C_\psi=3\sqrt3 e^{-3/2}+3e^{-1/2}$,
\begin{equation}
 \|D\Phi_J(x)\|_2\le L_\Phi=
 \max_j s_j\left(1+C_\psi^2\sum_{a=0}^{J-1}2^a4^a\right)^{1/2}.
\end{equation}
The bound follows by separately maximizing $|t|^3e^{-t^2/2}$ and
$3|t|e^{-t^2/2}$. Integration on a line segment in the convex box establishes
$\|\Phi_J(x)-\Phi_J(y)\|_2\le L_\Phi\|x-y\|_2$.
Increasing resolution increases $L_\Phi$ and may worsen robustness; fine scales
are not automatically an improvement.

\subsection{Deep counted likelihood and explicit sensitivities}
For column-vector notation let $h_0=\Phi_J(x)$,
$h_l=\operatorname{ReLU}(W_lh_{l-1}+b_l)$, $l=1,\ldots,L$, and
$f_\theta(x)=\sigma(W_{L+1}h_L+b_{L+1})\in(0,1)^m$.
For development counts $k_{ni}$ out of $b_{ni}$ trials, fit
\begin{equation}
 \mathcal L(\theta)= -\frac{1}{\sum_{n,i}b_{ni}}
 \sum_{n,i}\{k_{ni}\log f_{ni}+(b_{ni}-k_{ni})\log(1-f_{ni})\}
 +\frac\lambda2\sum_l\|W_l\|_F^2.
\end{equation}
Binomial counts retain different trial denominators. Treating all observed
fractions as equally precise generally changes this objective. If the counts
are dependent, the expression is a working loss, not an independent binomial
likelihood. Training, early-stopping validation and final audits must use
separate groups; neither LLM proposals nor final certification trials provide
training labels.

Away from ReLU kinks the raw-input derivative is
\begin{equation}
 Df_\theta(x)=D_\sigma W_{L+1}D_LW_L\cdots D_1W_1D\Phi_J(x),
\end{equation}
where $D_l$ contains activation indicators and
$D_\sigma=\operatorname{diag}(f_i(1-f_i))$. The implementation stores transposed
weights because observations are rows. At ReLU kinks it uses the zero
subgradient. For logits clipped to $[-35,35]$, it additionally sets the
logit derivative to zero outside that interval. Thus it differentiates the
actual deployed predictor, rather than an unclipped surrogate.

Because ReLU and clipping are one-Lipschitz and $\sigma'\le1/4$,
\begin{equation}
 \|f_\theta(x)-f_\theta(y)\|_2\le
 L_f\|x-y\|_2,\quad L_f=\tfrac14 L_\Phi\prod_{l=1}^{L+1}\|W_l\|_2.
\end{equation}
This global bound is conservative; local Jacobians can be much smaller.
For coordinate $i$, the proposal margin derivative is
$D\widehat z_i=-Df_i/\tau_i$. A useful sensitivity report gives raw units,
feature ranges, point derivatives and the spectral bound. It never labels
$\widehat z_i$ an audited margin. If a separate uniform approximation guarantee
$\sup_x|r_i(x)-f_i(x)|\le\epsilon_i$ were established, the triangle inequality
would imply $r_i(y)\le f_i(x)+L_f\|y-x\|+\epsilon_i$. A held-out mean log loss
does not establish that uniform guarantee.

\subsection{A finite-family audit with a stated shift radius}
Fix the family, sample sizes and weights before final outcomes are observed.
For coordinate $i$, let $K_i\sim\operatorname{Bin}(n_i,p_i)$ under audit law
$P_i$, and let $Q_i$ be the intended deployment attack law. Assume an externally
justified total-variation bound
\begin{equation}
 \operatorname{TV}(P_i,Q_i)=\sup_E|P_i(E)-Q_i(E)|\le\rho_i.
\end{equation}
For weights $w_i>0$, $\sum_iw_i\le1$, use the exact one-sided
Clopper--Pearson bound
\begin{equation}
 U_i=\begin{cases}
 \operatorname{Beta}^{-1}(1-\alpha w_i;K_i+1,n_i-K_i),&K_i<n_i,\\
 1,&K_i=n_i.
 \end{cases}
 \quad U_i^Q=\min\{1,U_i+\rho_i\}.
\end{equation}
With probability at least $1-\alpha$, all deployment risks obey
$r_i(Q_i)\le U_i^Q$. To prove this, each binomial bound fails with probability
at most $\alpha w_i$; the union bound controls their simultaneous failure.
On the complementary event, the definition of total variation applied to the
attack-success event gives $r_i(Q_i)\le p_i+\rho_i\le U_i+\rho_i$.
No independence across coordinates is needed. Independent Bernoulli trials
within each coordinate, a fixed sample size and a fixed target law are needed.
Random data-dependent radii require their own simultaneous error budget.
A feature displacement and a model Lipschitz bound do not justify a TV radius.

The shifted audited position is $z_i^Q=1-U_i^Q/\tau_i$; all must be nonnegative.
This is a conditional statistical guarantee given the external shift model,
not proof that the radius is true or that unregistered harms are absent.
At zero successes, $U_i=1-(\alpha w_i)^{1/n_i}$, so passage requires
\begin{equation}
 n_i\ge\left\lceil\frac{\log(\alpha w_i)}
 {\log(1-\tau_i+\rho_i)}\right\rceil,\qquad 0\le\rho_i<\tau_i.
\end{equation}
If $\rho_i\ge\tau_i$, zero observed successes cannot supply a finite passing
budget. Repeatedly inspecting these fixed-size bounds and stopping at passage
invalidates their stated coverage; use the existing sequential gate instead.

\subsection{Audit cost and scaling}
For positive costs $c_i$, set $a_i=c_i/[-\log(1-\tau_i+\rho_i)]$.
Ignoring integer ceilings, minimize
$C(w)=\sum_i a_i\log(1/(\alpha w_i))$ over positive weights summing to one.
The Lagrangian derivative is $-a_i/w_i+\nu=0$; summing gives
\begin{equation}
 w_i^*=a_i/\sum_ja_j.
\end{equation}
The Hessian is diagonal with entries $a_i/w_i^2>0$, proving a unique minimum.
The integer plan rounds upward and checks the exact zero-success inequality.
It is a planning rule conditional on zero successes, not a promised audit cost
under positive risk and not an optimizer of the discontinuous integer objective.
Freeze weights from costs and declared tolerances before audit outcomes.

Feature construction costs $O(Nd2^J)$. A dense network forward pass costs
$O(N\sum_l d_{l-1}d_l)$; storing every raw-input sensitivity additionally costs
$O(Nmd)$ for the output alone. The implementation supports arbitrary $d$ but
limits $J$ to eight to avoid accidental dictionary growth. Dimension-free
statistical safety does not imply dimension-free computational or predictive
accuracy. LLM scenario injections continue through the existing review and
measured-development-evidence pipeline; modifying a dictionary or threat law
requires a new proposal-model version and new final audits.
